\begin{pastexam}[label={px:gk3-107-ds-02}]{高考三級資訊處理 107 年 第二題}
（申論；25 分）

一非空的二元樹（binary tree），如果有 n$_0$ 個葉節點（leaf node）且 n$_2$ 個節點之分支度（degree）為 2，請證明
n$_0$ = n$_2$+1。（25 分）
\end{pastexam}

\begin{pastsoln}
和歷屆試題~\ref{px:gk3-099-ds-03} 的 ㈢ 相同（資結書定理~7.2）。 那裡用「數邊」證明；這題配 25 分，下面先寫
同一個數邊的證明，再補一個對節點數做歸納的證明，兩個寫一個就夠。

設二元樹有 $n \ge 1$ 個節點，分支度為 $0, 1, 2$ 的節點各有 $n_0, n_1, n_2$ 個。

\textbf{證明一（數節點和邊）。} 每個節點的分支度是 $0$、$1$ 或 $2$，所以
\[
n = n_0 + n_1 + n_2 .
\]
除了根以外，每個節點都恰好有一條邊從父節點連下來，根沒有，所以邊數 $= n - 1$。 換從父節點那一端數：分支度
為 $1$ 的節點往下連 $1$ 條邊，分支度為 $2$ 的連 $2$ 條，葉節點不連，每條邊恰好被數一次，所以邊數
$= n_1 + 2n_2$。 兩種數法相等：
\[
n_0 + n_1 + n_2 - 1 = n_1 + 2n_2 \quad\Longrightarrow\quad n_0 = n_2 + 1. \qquad\blacksquare
\]

\textbf{證明二（對節點數 $n$ 歸納）。}
\begin{itemize}
\item \textbf{基底} $n = 1$：只有根，它是葉節點，$n_0 = 1$、$n_2 = 0$，$n_0 = n_2 + 1$ 成立。
\item \textbf{歸納步驟}：設 $n \ge 2$，並假設所有節點數小於 $n$ 的非空二元樹都滿足 $n_0 = n_2 + 1$。 取一棵
$n$ 個節點的二元樹 $T$，拿掉它的任一片葉子 $x$（非空有限樹一定有葉子），得到 $n - 1$ 個節點的非空二元樹
$T'$。 令 $p$ 是 $x$ 的父節點（$n \ge 2$，所以 $x$ 不是根，$p$ 存在）。
\begin{itemize}
\item 若 $p$ 在 $T$ 中分支度為 $1$：拿掉 $x$ 後 $p$ 變成葉子。 葉子數不變（少了 $x$、多了 $p$），分支度 $2$
的節點數不變。 所以 $n_0(T) = n_0(T')$、$n_2(T) = n_2(T')$。
\item 若 $p$ 在 $T$ 中分支度為 $2$：拿掉 $x$ 後 $p$ 分支度變成 $1$。 葉子少 $1$，分支度 $2$ 的節點也少 $1$：
$n_0(T) = n_0(T') + 1$、$n_2(T) = n_2(T') + 1$。
\end{itemize}
兩種情況下 $n_0(T) - n_2(T) = n_0(T') - n_2(T')$，由歸納假設等於 $1$，所以 $n_0(T) = n_2(T) + 1$。
\end{itemize}
由數學歸納法，所有非空二元樹都滿足 $n_0 = n_2 + 1$。 \hfill$\blacksquare$

\textbf{注意}：題目說「非空」是必要的，空樹的 $n_0 = n_2 = 0$，式子不成立。 證明一用到「邊數 $= n - 1$」，
這要 $n \ge 1$ 才對。 例子：三個節點的完整二元樹 $n_0 = 2$、$n_2 = 1$；一條三個節點的鏈 $n_0 = 1$、$n_2 = 0$，
都符合。
\end{pastsoln}
